% ============================================================================
\chapter{서론}\label{ch:intro}
% ============================================================================

\section{연구 배경}\label{sec:background}

갑상선 결절은 매우 흔한 소견이다. 고해상도 초음파를 사용하면 무작위로 선정한 성인의 19--68\%에서
결절이 발견되며, 그중 임상적으로 문제가 되는 갑상선암은 7--15\% 수준으로 보고된다\citep{haugen2016ata}.
발견되는 결절의 대부분이 양성이므로, 진료의 핵심은 어떤 결절에 세침흡인검사가 필요한지를
가려내는 데 있다. 이 판단의 1차 근거가 초음파 영상이다.

미국영상의학회의 ACR TI-RADS를 비롯한 주요 판독 체계는 결절의 구성, 에코, 형태, 경계(margin),
에코 발생 병소를 각각 채점하여 합산한다\citep{tessler2017tirads}. 이 가운데 경계는 독립된 채점
항목으로, 매끄럽거나 불분명한 경계에는 0점, 분엽상 또는 불규칙한 경계에는 2점,
갑상선 외 침범에는 3점이 부여된다. 경계 점수 하나가 결절의 위험 등급을 바꾸고 결과적으로
세침흡인검사 여부를 좌우할 수 있다. 다시 말해 갑상선 결절의 자동 분할에서 경계는 부수적인
정보가 아니라 판독 결과에 직접 들어가는 정보이다.

그러나 초음파 영상에서 경계를 정확히 그리는 일은 사람에게도 어렵다. 스페클 잡음, 결절 하방의
음향 음영, 주변 조직과의 낮은 대조도가 경계를 흐리게 만들고, 경계 관련 소견은 판독자 간
일치도가 낮은 항목으로 알려져 있다\citep{hoang2018interobserver}. 딥러닝 기반 자동 분할은
이 변이를 줄이고 판독 시간을 단축할 수단으로 연구되어 왔다\citep{ronneberger2015unet,gong2021multitask,gong2023trfe}.

자동 분할이 실제 진료에 들어가려면 정확도와 함께 배포 조건도 충족해야 한다. 건강검진 현장이나
휴대형 초음파 장비, 소규모 의원의 판독 PC에는 GPU가 없는 경우가 많다. 고성능 GPU 서버를
전제로 설계된 모델은 이런 환경에서 실시간으로 동작하지 못한다. 따라서 경계 정확도를 높이는
기법이 추론 비용을 늘린다면, 그 기법은 바로 이 배포 조건과 충돌하게 된다.

\section{문제 정의}\label{sec:problem}

본 연구는 위 배경에서 다음 세 가지 문제를 정의한다. 각 문제는 이후 실험 설계의 단계에 그대로
대응된다.

\paragraph{P1. 경계 품질은 영역 지표에 가려진다.}
분할 모델의 성능은 흔히 mAP나 IoU 같은 영역 중심 지표로 보고된다. 이 지표들은 결절 전체 면적에
대한 겹침을 재기 때문에, 면적에서 차지하는 비율이 작은 경계 부근의 몇 픽셀 오차에는 둔감하다.
경계가 실제로 나아졌는지 판정하려면 Boundary IoU\citep{cheng2021boundaryiou}나 95 백분위
Hausdorff 거리(HD95)\citep{taha2015metrics}처럼 경계에 특화된 지표로 함께 평가해야 한다.

\paragraph{P2. 공개 데이터셋은 서로 이질적이다.}
공개된 갑상선 초음파 데이터셋은 수집 기관, 장비, 화각, 주석 방식이 서로 다르다. 한 데이터셋에서
학습한 모델이 다른 데이터셋에서도 같은 성능을 낼지는 알 수 없으며, 의료영상 전반에서 이러한
도메인 이동은 성능 저하의 주된 원인으로 지적된다\citep{guan2022domain}. 여러 데이터셋을 혼합해
학습하고, 학습에 쓰지 않은 데이터셋에서 직접 평가하는 실험이 필요하다.

\paragraph{P3. 배포 환경에는 GPU가 없다.}
GPU 없이 CPU만으로 실시간 추론이 가능한 수준까지 모델의 크기와 연산량을 줄여야 한다.
이때 경량화가 P1에서 얻은 경계 품질을 훼손하지 않는지도 함께 확인해야 한다.

\section{연구 목적 및 가설}\label{sec:hypothesis}

본 연구의 목적은 갑상선 결절 초음파 영상에서 경계 정확도를 높이면서도 CPU에서 실시간으로
동작하는 검출·분할 모델을 만드는 것이다. 이를 위해 다음 세 가설을 세우고, 각 가설마다 검증 방법과
기각 조건을 미리 정하였다.

\begin{description}[leftmargin=2.2em,style=nextline]
  \item[H1 (모델 용량 가설)]
    YOLOv11-seg nano는 small보다 파라미터가 약 1/3.5임에도 갑상선 결절 검출·분할에서
    동등한 정확도를 낸다.\\
    \textit{검증}: 동일한 5-fold 분할과 학습 설정에서 두 모델을 학습하고 held-out 집합의
    Mask \mAP{}를 비교한다. \textit{기각 조건}: nano가 small보다 Mask \mAP{}에서
    fold 표준편차 이상 뒤처지는 경우.

  \item[H2 (경계 정교화 가설)]
    학습 시에만 동작하는 보조 경계 분기와 경계 가중 손실은 추론 비용을 전혀 늘리지 않고
    경계 품질을 개선한다.\\
    \textit{검증}: Dice, Boundary IoU, HD95를 영상 단위 대응표본 검정으로 비교하고, 배포 그래프의
    GFLOPs가 베이스라인과 같은지 확인한다. \textit{기각 조건}: 경계 지표의 개선이 유의하지 않거나
    추론 연산량이 증가하는 경우.

  \item[H3 (경량화 가설)]
    INT8 사후 양자화로 CPU 배포 요건을 만족시키면서 정확도 손실은 무시할 수 있는 수준으로
    유지된다.\\
    \textit{검증}: PyTorch, ONNX FP32·FP16, ONNX Runtime INT8, OpenVINO INT8의 다섯 정밀도에서
    정확도·크기·지연 시간을 비교한다. \textit{기각 조건}: CPU 지연이 실시간 기준에 못 미치거나
    Dice 손실이 fold 표준편차를 넘는 경우.
\end{description}

\section{연구의 기여}\label{sec:contribution}

본 연구의 기여는 다음과 같다.

\begin{enumerate}
  \item \textbf{추론 비용이 없는 경계 정교화.} 분할 헤드의 프로토타입 특징에서 갈라지는 학습 전용
    경계 분기와 형태학적 경계 밴드 기반 가중 손실을 제안하였다. 추가 파라미터는 4,657개이고
    추론 그래프는 베이스라인과 동일하며, held-out 평가에서 HD95를 13.0\% 줄였다.
  \item \textbf{경계 특화 평가와 통계 검정.} mAP에 드러나지 않는 경계 개선을 Boundary IoU와 HD95로
    측정하고, 평균 비교에 그치지 않고 영상 단위 대응표본 검정, 부트스트랩 신뢰구간,
    효과크기로 개선이 우연의 범위를 벗어나는지 확인하였다.
  \item \textbf{교차 데이터셋 결과의 신중한 보고.} 교차 데이터셋 실험을 3개 시드로 반복하여,
    단일 실행에서 크게 악화된 것처럼 보이던 차이가 실행 간 변동 범위 안에 있음을 보였다.
    또한 두 데이터셋의 정량적 차이(결절 수 분포, 결절 크기, 해상도)를 분석하여 교차 데이터셋
    성능이 낮은 원인을 데이터 측면에서 설명하였다.
  \item \textbf{배포 형식별 경량화 비교.} 같은 INT8이라도 변환 도구에 따라 CPU 지연이 44.00 ms와
    4.10 ms로 10배 이상 차이 남을 보이고, 그 원인을 양자화에서 제외되는 연산의 범위로 설명하였다.
    구현 과정에서 확인한 실패 사례(검출 헤드 양자화 시 출력 붕괴 등)도 함께 기록하였다.
\end{enumerate}

\section{논문의 구성}\label{sec:organization}

본 논문의 구성은 다음과 같다. \chapref{ch:related}에서는 의료영상 분할, 실시간 인스턴스 분할,
경계 정교화, 모델 경량화, 분할 평가 지표에 관한 선행연구를 정리하고 연구 공백을 도출한다.
\chapref{ch:dataset}에서는 사용한 두 데이터셋의 특성과 정량적 차이, 전처리와 분할 설계를 설명한다.
\chapref{ch:method}에서는 제안하는 경계 정교화 모듈과 경계 인식 손실, 학습--추론 분리 구조,
양자화 파이프라인을 기술한다. \chapref{ch:experiments}에서는 실험 설정과 결과를 가설 순서에 따라
제시하고, \chapref{ch:discussion}에서는 결과의 의미와 연구의 한계를 논의한다.
마지막으로 \chapref{ch:conclusion}에서 결론과 향후 연구 방향을 제시한다.
